\documentclass[14pt]{extarticle}
\input{../../../scripts/tex_preamble.tex}

\title{ورقة تمرن 2 للصف العاشر 10}

\begin{document}

\maketitle
\thispagestyle{fancy}

\begin{enumerate}[itemsep=2cm]


\item
في مطعم الطويل ثمن ساندويتش فلافل 13 شاقل وثمن ساندويتش شاورما 22 شاقل. اكتب برنامجًا يستقبل عدد ساندويتشات الفلافل وعدد ساندويتشات الشاورما التي باعها المطعم في أحد الأيام. على البرنامج أن يحسب ويطبع المدخول الكلي الذي جمعه المطعم في ذلك اليوم.
{\small
\ifdetailed
\begin{boxExample}
\begin{english}
\begin{minted}{text}
Enter falafel sandwiches:
> 35
Enter shawarma sandwiches:
> 46
Total income: 1467
\end{minted}
\end{english}
\end{boxExample}
\fi
\ifwithsols
\begin{boxSolution}
\begin{english}
\begin{minted}{csharp}
Console.WriteLine("Enter falafel sandwiches:");
int falafel = int.Parse(Console.ReadLine());
Console.WriteLine("Enter shawarma sandwiches:");
int shawarma = int.Parse(Console.ReadLine());
int total = falafel * 13 + shawarma * 22;
Console.WriteLine("Total income: " + total);
\end{minted}
\end{english}
\end{boxSolution}
\clearpage
\fi
}

\item
في محل لبيع الاجهزة الكهربائية يتم دفع مبلغ الشراء بالطريقة التالية: نصف المبلغ نقدي والنصف الآخر يتم تقسيمه على خمسة أقساط. اكتب برنامجًا يستقبل الثمن الكلي للأجهزة الكهربائية التي تم شراؤها، على البرنامج أن يحسب ويطبع المبلغ المدفوع نقدا وقيمة القسط الواحد.
{\small
\ifdetailed
\begin{boxExample}
\begin{english}
\begin{minted}{text}
Enter the total price:
> 10000
Cash payment: 5000
Installment payment: 1000
\end{minted}
\end{english}
\end{boxExample}
\fi
\ifwithsols
\begin{boxSolution}
\begin{english}
\begin{minted}{csharp}
Console.WriteLine("Enter the total price:");
double total = double.Parse(Console.ReadLine());
double cash = total / 2;
double installment = cash / 5;
Console.WriteLine("Cash payment: " + cash);
Console.WriteLine("Installment payment: " + installment);
\end{minted}
\end{english}
\end{boxSolution}
\fi
}
\clearpage

\item
اكتب برنامجًا يستقبل علامة الفصل الأول وعلامة الفصل الثاني وعلامة الفصل الثالث بموضوع علم الحاسوب، على البرنامج أن يحسب ويطبع معدل العلامات بالطريقة التالية: 30\% من علامة الفصل الأول، 25\% من علامة الفصل الثاني و 45\% من علامة الفصل الثالث.
{\small
\ifdetailed
\begin{boxExample}
\begin{english}
\begin{minted}{text}
Enter grade 1:
> 82
Enter grade 2:
> 77
Enter grade 3:
> 92
Final grade: 85.25
\end{minted}
\end{english}
\end{boxExample}
\fi
\ifwithsols
\begin{boxSolution}
\begin{english}
\begin{minted}{csharp}
Console.WriteLine("Enter grade 1:");
int g1 = int.Parse(Console.ReadLine());
Console.WriteLine("Enter grade 2:");
int g2 = int.Parse(Console.ReadLine());
Console.WriteLine("Enter grade 3:");
int g3 = int.Parse(Console.ReadLine());
double finalGrade = g1 * 0.30 + g2 * 0.25 + g3 * 0.45;
Console.WriteLine("Final grade: " + finalGrade);
\end{minted}
\end{english}
\end{boxSolution}
\clearpage
\fi
}


\item
اكتب برنامجًا يستقبل من المستخدم الميزانية لمشروع ما. \\
ثم يستقبل من المستخدم 3 مصروفات للمشروع. \\
بعد كل مرة يدخل فيها المستخدم مصروفًا، اطبع له الميزانية المتبقية.
{\small
\ifdetailed
\begin{boxExample}
\begin{english}
\begin{minted}{text}
Enter the budget:
> 1000
Enter expense 1:
> 200
Remaining budget: 800
Enter expense 2:
> 150
Remaining budget: 650
Enter expense 3:
> 50
Remaining budget: 600
\end{minted}
\end{english}
\end{boxExample}

\ifwithsols
\begin{boxSolution}
\begin{english}
\begin{minted}{csharp}
int budget;
int expenses;
Console.Write("Please enter the budget:");
budget = int.Parse(Console.ReadLine());
Console.WriteLine("Please enter the first expense:");
expenses = int.Parse(Console.ReadLine());
budget = budget - expenses;
Console.WriteLine(budget);

Console.WriteLine("Please enter the second expense:");
expenses = int.Parse(Console.ReadLine());
budget = budget - expenses;
Console.WriteLine(budget);

Console.WriteLine("Please enter the third expense:");
expenses = int.Parse(Console.ReadLine());
budget = budget - expenses;
Console.WriteLine(budget);
\end{minted}
\end{english}
\end{boxSolution}
\fi

\fi
}

\clearpage

\item
اكتب برنامجًا يستقبل من المستخدم السعر الأولي لمنتج ما. \\
ثم يستقبل من المستخدم 3 تخفيضات على هذا المنتج بالنسبة المئوية. \\
بعد كل مرة يدخل فيها المستخدم تخفيضًا، اطبع له السعر الجديد للمنتج.
{\small
\ifdetailed
\begin{boxExample}
\begin{english}
\begin{minted}{text}
Enter the initial price:
> 500
Enter discount 1 (%):
> 20
Price after discount:400
Enter discount 2 (%):
> 10
Price after discount: 360
Enter discount 3 (%):
> 5
Price after discount: 342
\end{minted}
\end{english}
\end{boxExample}

\ifwithsols
\begin{boxSolution}
\begin{english}
\begin{minted}{csharp}
double price;
int discount;
Console.WriteLine("Please enter the initial price:");
price = double.Parse(Console.ReadLine());
Console.WriteLine("Please enter the first discount (%):");
discount = int.Parse(Console.ReadLine());
price = price * (100 - discount) / 100;
Console.WriteLine("The new price = " + price);

Console.WriteLine("Please enter the second discount (%):");
discount = int.Parse(Console.ReadLine());
price = price * (100 - discount) / 100;
Console.WriteLine("The new price = " + price);

Console.WriteLine("Please enter the third discount (%):");
discount = int.Parse(Console.ReadLine());
price = price * (100 - discount) / 100;
Console.WriteLine("The final price = " + price);
\end{minted}
\end{english}
\end{boxSolution}
\clearpage
\fi

\fi
}

\item
اكتب برنامجًا يبدأ بالحساب من 3600 ثانية (أي ساعة واحدة)، ثم يستقبل من المستخدم عدد الثواني التي مرت. \\
يستقبل ذلك 4 مرات، وفي كل مرة يطبع الوقت المتبقي بالثواني.
{\small
\ifdetailed
\begin{boxExample}
\begin{english}
\begin{minted}{text}
Enter seconds passed:
> 600
Remaining time (seconds): 3000
Enter seconds passed:
> 300
Remaining time (seconds): 2700
Enter seconds passed:
> 1200
Remaining time (seconds): 1500
Enter seconds passed:
> 500
Remaining time (seconds): 1000
\end{minted}
\end{english}
\end{boxExample}

\ifwithsols
\begin{boxSolution}
\begin{english}
\begin{minted}{csharp}
int time = 3600;
int passedTime;
Console.WriteLine("Enter seconds passed:");
passedTime = int.Parse(Console.ReadLine());
time = time - passedTime;
Console.WriteLine($"Remaining time (seconds): {time}");

Console.WriteLine("Enter seconds passed:");
passedTime = int.Parse(Console.ReadLine());
time = time - passedTime;
Console.WriteLine($"Remaining time (seconds): {time}");

Console.WriteLine("Enter seconds passed:");
passedTime = int.Parse(Console.ReadLine());
time = time - passedTime;
Console.WriteLine($"Remaining time (seconds): {time}");

Console.WriteLine("Enter seconds passed:");
passedTime = int.Parse(Console.ReadLine());
time = time - passedTime;
Console.WriteLine($"Remaining time (seconds): {time}");
\end{minted}
\end{english}
\end{boxSolution}
\fi
\fi
}

\clearpage

\item
اكتب برنامجًا يستقبل عددين صحيحين وموجبين ثنائيي المنزلة، على البرنامج ان يطبع حاصل ضرب خانتي الآحاد من العددين.
{\small
\ifdetailed
\begin{boxExample}
\begin{english}
\begin{minted}{text}
Enter first number:
> 15
Enter second number:
> 23
Units product: 15
\end{minted}
\end{english}
\end{boxExample}
\fi
\ifwithsols
\begin{boxSolution}
\begin{english}
\begin{minted}{csharp}
Console.WriteLine("Enter first number:");
int num1 = int.Parse(Console.ReadLine());
Console.WriteLine("Enter second number:");
int num2 = int.Parse(Console.ReadLine());
int units1 = num1 % 10;
int units2 = num2 % 10;
int product = units1 * units2;
Console.WriteLine("Units product: " + product);
\end{minted}
\end{english}
\end{boxSolution}
\clearpage
\fi
}

\item
في مصنع لإنتاج حبات الشوكولاتة يعبئون كل 30 حبة بعلبة. اكتب برنامجًا يستقبل عدد حبات الشوكولاتة التي أنتجها المصنع في أحد الايام، على البرنامج أن يطبع عدد العُلب التي ملأها المصنع ذلك اليوم، وأيضًا يطبع عدد حبّات الشوكولاتة التي بقيت ولم تكفِ لملء علبة إضافية.
{\small
\ifdetailed
\begin{boxExample}
\begin{english}
\begin{minted}{text}
Enter the number of chocolates:
> 3070
Full boxes: 102
Leftover chocolates: 10
\end{minted}
\end{english}
\end{boxExample}
\fi
\ifwithsols
\begin{boxSolution}
\begin{english}
\begin{minted}{csharp}
Console.WriteLine("Enter the number of chocolates:");
int chocolates = int.Parse(Console.ReadLine());
int boxes = chocolates / 30;
int leftovers = chocolates % 30;
Console.WriteLine("Full boxes: " + boxes);
Console.WriteLine("Leftover chocolates: " + leftovers);
\end{minted}
\end{english}
\end{boxSolution}
\fi
}



\end{enumerate}

\vspace{3cm}
\begin{flushleft}
أرجو لكم وقتًا ممتعًا.

الأستاذ محمود اغبارية.
\end{flushleft}

\end{document}
